\chapter{Literature Survey}

This chapter looks at the work that other researchers have already done in the areas this
project touches. It covers wearable EEG devices, publicly available motor imagery
datasets, the methods used to extract features and classify EEG signals, and the use of
EEG to measure a person's mental and cognitive state. A short section on decoding speech
from EEG is also included, since that is the direction this work is expected to move in
later. The aim of this chapter is to understand what has already been done, what problems
are still open, and where this project fits in.

%------------------------------------------------------------------
\section{Wearable EEG Systems}
%------------------------------------------------------------------

EEG devices have come a long way from the large machines used in hospitals. Today there
are small, portable EEG headsets that a person can wear while going about normal
activities. This matters for any brain-computer interface (BCI) that is meant to be used
outside a lab.

A real-time system for neuroimaging and cognitive monitoring using wearable dry EEG was
built and tested \cite{mullen2015wearable}. Dry electrodes do not need any gel, which
makes them much faster to set up and more comfortable to wear for long periods. The system
could track brain activity as it happened, which is useful for practical applications. The
results showed that wearable dry EEG gives reasonably good signal quality, though it is
not as clean as gel-based EEG recorded in a lab.

A different approach was explored by placing dry-contact electrodes inside the ear to
record EEG \cite{kappel2019earEEG}. This is called ear-EEG and it is convenient because
the device looks like a regular earphone. The study tested whether ear-mounted electrodes
could pick up brain signals well enough to be useful. The signals were weaker than those
from standard head-mounted EEG, but the setup was far less intrusive, which makes it more
suitable for daily use.

Both of these studies are directly relevant to the hardware used in this project. The GELO
BCI headset used here is a dry-electrode headphone with only four EEG channels (A1, A2, C3
and C4) and a single reference at Cz, sampled at 250 Hz. It gains the comfort and quick
setup reported in these studies, but it also inherits their limitations: fewer channels,
weaker signals and more noise than a research-grade cap. This trade-off shapes much of the
work reported in the later chapters.

%------------------------------------------------------------------
\section{Motor Imagery EEG Datasets}
%------------------------------------------------------------------

Motor imagery (MI) means imagining a movement without actually making it. When a person
imagines moving the left or right hand, the activity over the motor cortex changes, and
these changes can be picked up by electrodes placed at C3 and C4. Building a classifier for
this needs a good amount of labelled data, which is hard for a student project to collect
from scratch.

A large open EEG dataset for left- and right-hand motor imagery in acute stroke patients
was published and is used as the reference dataset in this project \cite{liu2024mi}. The
dataset contains recordings from 50 patients between 31 and 77 years of age, taken 1 to 30
days after stroke. Each patient performed 40 trials, 20 imagining a left-hand grasp and 20
imagining a right-hand grasp, giving 2000 trials in total. Each trial followed a fixed
pattern of a 2 second instruction, 4 seconds of motor imagery and a 2 second break. The
data was recorded with a wireless portable saline EEG system at 500 Hz using 30 EEG
channels, 2 EOG channels and one marker channel, with CPz as the reference. Both raw data
and preprocessed data are provided, along with the clinical details of each patient.

This dataset is important for the present work for two reasons. First, it supplies a
labelled, well-documented set of left and right MI trials that can be used to train a
classifier, which is the first objective of this project. Second, the paper reports its own
analysis of the data, including event-related desynchronisation and a set of classification
results, which gives a clear baseline to compare against. The channel layout is also
convenient, as C3 and C4 are available in the dataset and are the same two positions
recorded by the headset used here. The main limitation is that the subjects are stroke
patients recorded with a 30-channel system, while the data collected in this project comes
from healthy subjects using only four channels. This difference between the training data
and the test data turns out to matter a great deal, as the results chapter shows.

%------------------------------------------------------------------
\section{Feature Extraction and Classification for Motor Imagery}
%------------------------------------------------------------------

Once the data is available, the next problem is turning the raw signal into features and
training a model on them. This is harder than it sounds, because EEG signals vary a lot
between people and even between sessions for the same person.

The standard way of looking at motor imagery is through event-related desynchronisation
(ERD) and event-related synchronisation (ERS). ERD is a drop in power in a given frequency
band during a task and ERS is a rise in power. For motor imagery, power in the alpha
(8--15 Hz) and beta (15--30 Hz) bands usually drops over the motor cortex on the side
opposite to the imagined hand \cite{pfurtscheller1999erd}. The reference dataset paper computes this by filtering C3
and C4 between 8 and 30 Hz, squaring the signal to get power, averaging across trials and
comparing against a baseline taken just before the task \cite{liu2024mi}. The same paper
also benchmarks several classifiers on its own data for the left versus right MI problem,
using the 8--30 Hz band and a 4 second window after the cue. Common Spatial Pattern with
Linear Discriminant Analysis (CSP+LDA) \cite{ramoser2000csp} reached about 55.6\%, Filter
Bank CSP with a Support Vector Machine (FBCSP+SVM) \cite{cortes1995svm} about 57.6\%, a Riemannian approach about 61.2\%, and the best
method reported, a temporal window filter bank combined with a Riemannian mean classifier,
reached 72.21\%. These numbers are a useful reminder that even with 30 channels and clean
laboratory recordings, binary motor imagery classification is a difficult problem, and that
accuracies in the 50--70\% range are normal rather than poor.

A review of deep learning methods for EEG classification and feature extraction gives a
broader picture of what is available beyond CSP and SVM \cite{xu2023review}. It goes
through the common network types used on EEG, explains how the signal is usually prepared
before it reaches the network, and points out the main difficulty, which is that EEG
datasets are small compared with the amount of data deep networks normally need. This is
why a simple classifier such as an SVM is a sensible starting point for a project working
with a small, self-recorded dataset.

A more recent model combines a convolutional network with transformer self-attention for
EEG decoding \cite{song2023conformer}. The convolutional part picks up short-term patterns
in the signal while the attention part links information across the whole trial, and the
model also produces visualisations that show which parts of the signal it is relying on.
It reports better accuracy than CSP-based methods on standard motor imagery datasets. This
is a clear option for improving the classifier in future work, although it needs more
training data and more computation than the present setup provides.

Transfer learning has been used to improve imagined speech recognition across different
users \cite{jerrin2021imagined}. Here a model is first trained on data from many people and
then adjusted for a new person using only a small amount of their data. This cut down the
time needed to set the system up for each new user and improved how well the model worked
across subjects. Even so, accuracy still dropped noticeably from one person to another,
which shows that handling differences between individuals is still an open problem. The
same problem appears in this project, where a model trained on one dataset is tested on
data recorded from a different group of people with a different device.

An adaptive hybrid brain-computer interface for helping stroke patients recover hand
function has also been proposed \cite{su2024hybrid}. The system combined motor imagery with
steady-state visual evoked potentials and used adaptive learning, meaning the model updated
itself as the user's brain patterns changed over time. Combining two kinds of signal raised
the classification accuracy, and the adaptive part helped the model keep up with changes in
the user. The reported improvements in rehabilitation outcomes show that a well-trained EEG
model can have real clinical value.

%------------------------------------------------------------------
\section{Decoding Speech from EEG}
%------------------------------------------------------------------

A related line of research tries to work out what a person is saying, or only thinking of
saying, from EEG alone. This is not part of the current objectives but it is the direction
this project is expected to extend into, so the main findings are worth noting.

EEG was recorded from people while they were speaking, and it was shown that the brain
signals contain patterns matching what the person was saying, even when no sound was
produced or the speech was noisy \cite{krishna2019nospeech}. This mattered because it meant
EEG on its own carries enough information about speech intent to be used for recognition.
The experiments were done in a controlled lab, so how well this holds in noisy everyday
conditions is still not known. Later work went further and tried to rebuild speech from EEG
signals, and found that EEG carries some information about how sounds are formed, which
allowed partial reconstruction of spoken words from brain activity alone
\cite{krishna2020eeg}. The method depended on specific tasks and had to be trained
separately for each person, which makes it hard to use on a new user without starting over.

EEG recorded during both imagined and spoken speech has also been used to classify the
types of sounds being thought about or said \cite{zhao2015phonological}. This confirmed
that imagined speech, where no sound is made at all, still leaves a pattern in EEG that can
be detected, which is the basic assumption behind any silent communication system. On the
listening side, the auditory EEG decoding challenge for ICASSP 2024 focused on working out
what a person is listening to from their EEG \cite{bollens2025auditory}. The dataset and the
evaluation method were released publicly, which helps the wider research community compare
results fairly.

%------------------------------------------------------------------
\section{EEG for Cognitive and Mental State Assessment}
%------------------------------------------------------------------

Apart from decoding what a person intends to do or say, EEG is also widely used to measure
how a person's mind is working at a given moment. This is the third objective of this
project, which is to compute a cognitive score from EEG frequency bands.

Mental state classification using an EEG-based brain-machine interface has been studied by
extracting features from EEG and using machine learning to tell different cognitive
conditions apart \cite{bird2018mental}. The results showed that EEG can separate distinct
mental states with reasonable accuracy. The work looked at states in isolation, however,
and did not follow how a changing cognitive state affects performance on an actual task.

Deep learning has been applied to raw EEG to identify visual learners in real time
\cite{jawed2024visual}. Recordings from 34 subjects at rest and during learning tasks were
fed to LSTM, LSTM-CNN and LSTM-FCNN models, and the LSTM-CNN model reached about 94\%
accuracy. What is useful here is that the models worked on raw EEG without heavy offline
cleaning, which shows that real-time assessment from a continuous EEG stream is practical.

Drowsiness during driving has been estimated from EEG using deep Q-learning
\cite{ming2021drowsiness}, which treats the problem as one of learning from feedback rather
than fixed labels. This is a useful idea because drowsiness builds up gradually and does not
split neatly into classes. A more conventional system for the same problem used spectral
entropy and power spectral density as features, which describe how the brain's electrical
activity is spread across frequency bands, and fed them to boosting classifiers to detect
signs of fatigue \cite{jiri2025drowsiness}. This second approach is closest to the method
used in this project, where relative power in the delta, theta, alpha, beta and gamma bands
is computed from the PSD and then combined into scores. In particular, the stress score used
here is based on the ratio of relative high-beta power to relative alpha power, and the
fatigue score on the ratio of combined theta and alpha power to beta power, both of which
rest on the same band-power relationships these studies rely on. The limitation of such
systems is that they are tuned for one specific state, so the thresholds and reference
values have to be worked out again for any new measure.

%------------------------------------------------------------------
\section{Summary of the Literature Survey}
%------------------------------------------------------------------

The studies reviewed above map onto the three objectives of this project as follows.

\textbf{Building an EEG database for motor imagery.} Wearable dry-electrode systems,
including ear-mounted ones, can record usable EEG with a small number of channels and very
little setup \cite{mullen2015wearable, kappel2019earEEG}, and the open stroke-patient motor
imagery dataset provides a well-documented left and right hand MI recording protocol that
can be followed when collecting new data \cite{liu2024mi}.

\textbf{Designing a motor imagery classifier.} ERD and ERS analysis of C3 and C4 in the
8--30 Hz band is the accepted way to look at motor imagery, and CSP, filter bank CSP,
LDA, SVM and Riemannian methods form the standard set of classifiers, with reported
accuracies for left versus right MI ranging from about 55\% to 72\%
\cite{liu2024mi}. Deep learning offers further gains \cite{xu2023review,
song2023conformer}, transfer learning helps when moving between users
\cite{jerrin2021imagined}, and adaptive models help when a user's signals drift over time
\cite{su2024hybrid}.

\textbf{Computing a cognitive score from frequency bands.} Relative power in the EEG
frequency bands, together with PSD and entropy features, can be used to separate mental
states and to track attention and fatigue \cite{bird2018mental, jawed2024visual,
ming2021drowsiness, jiri2025drowsiness}.

Three gaps stand out from this review. First, almost all of the motor imagery results
reported above come from research-grade systems with 20 or more channels, while a practical
wearable device may have only a few; how much accuracy is lost in that reduction is not
well documented. Second, models are usually trained and tested on the same dataset and the
same recording device, so very little is known about what happens when a model trained on
one dataset is applied to data recorded elsewhere with different hardware. Third, work on
motor imagery decoding and work on cognitive state monitoring have largely been carried out
separately, even though both are computed from the same recording. This project addresses
these three points by recording left and right motor imagery on a four-channel dry
headset, training classifiers on the public dataset and testing them on the recorded data,
and computing a cognitive score from the same signal within a single system.